\documentclass[11pt, aspectratio=169]{beamer}
\usetheme{metropolis}
\usepackage{amsmath}

\setbeamerfont{title}{size=\Large, series=\bfseries}
\setbeamerfont{author}{size=\small}
\setbeamerfont{subtitle}{size=\normalsize}

\title{SonicCraft}
\subtitle{An Interactive Python Desktop Suite for Signal Analysis}
\author{Md Mehrab Hossain (2305108) \& Mahbubul Islam Mahi (2305115)}
\date{27 September 2026}

\begin{document}

\begin{frame}[plain]
  \titlepage
\end{frame}

% ==========================================
% PART 1: MEHRAB
% ==========================================
\begin{frame}[standout]
  Part 1: System Architecture \& Time-Domain \\
  \vspace{0.5em}
  \large Presenter: Md Mehrab Hossain \\
  \normalsize ID: 2305108
\end{frame}

\begin{frame}{Project Overview}
  \begin{itemize}[<+->]
    \item \textbf{Core Philosophy:} SonicCraft bridges foundational Signals \& Systems theory with a real-time, responsive desktop interface.
    \item \textbf{Decoupled Architecture:} The PyQt6 GUI presentation layer is strictly separated from the non-blocking NumPy/SciPy DSP execution core.
    \item \textbf{Presentation Flow:} We will cover Time-Domain operations first, followed by advanced Frequency-Domain transformations and spectral editing.
  \end{itemize}
\end{frame}

\begin{frame}{Time-Domain Operations \& LTI Systems}
  \begin{columns}[T]
    \begin{column}{0.5\textwidth}
      \textbf{Discrete Signal Manipulation}
      \begin{itemize}[<+->]
        \item Audio is loaded as discrete-time arrays via SoundFile.
        \item Amplitude scaling and envelope fading via element-wise array multiplication.
      \end{itemize}
    \end{column}
    
    \begin{column}{0.5\textwidth}
      \textbf{Digital LTI Filtering}
      \begin{itemize}[<+->]
        \item 9-Band Parametric EQ and standard IIR/FIR filters (SciPy).
        \item Represented by the linear difference equation:
      \end{itemize}
      \vspace{0.5em}
      \uncover<3->{
        \[ \sum_{k=0}^{N} a_k y[n-k] = \sum_{m=0}^{M} b_m x[n-m] \]
      }
    \end{column}
  \end{columns}
\end{frame}

% ==========================================
% PART 2: MAHI
% ==========================================
\begin{frame}[standout]
  Part 2: Frequency-Domain \& Spectral Editing \\
  \vspace{0.5em}
  \large Presenter: Mahbubul Islam Mahi \\
  \normalsize ID: 2305115
\end{frame}

\begin{frame}{Frequency-Domain Analysis}
  \textbf{Fast Fourier Transform (FFT)} \\
  Converts continuous-time theory into discrete computational bins to identify dominant frequencies in real-time.
  \[ X[k] = \sum_{n=0}^{N-1} x[n] e^{-j 2\pi k n / N} \]
  
  \pause
  \vspace{0.5em}
  \textbf{Short-Time Fourier Transform (STFT)} \\
  \begin{itemize}[<+->]
    \item Chops audio into windowed, overlapping blocks to create a 2D matrix (Time vs. Frequency vs. Magnitude).
    \item Visualized in the UI as an interactive Spectrogram heatmap.
  \end{itemize}
\end{frame}

\begin{frame}{Advanced DSP \& Spectral Editing}
  \begin{itemize}[<+->]
    \item \textbf{Spectral Masking:} Users can draw a bounding box over specific time/frequency coordinates on the Spectrogram to mathematically zero-out unwanted frequencies.
    \item \textbf{Image-to-Audio Reconstruction:} Uses PIL to read grayscale images, map pixels to frequency bins, and apply the Griffin-Lim algorithm to synthesize playable audio arrays.
    \item \textbf{Inverse STFT:} Both processes utilize the Inverse STFT to rebuild modified magnitude spectra back into a 1D time-domain waveform.
  \end{itemize}
\end{frame}

\begin{frame}{Interactive Workflow (Demo)}
  \begin{center}
    % UNCOMMENT AND ADD YOUR SCREENSHOT HERE
    % \includegraphics[width=0.8\textwidth]{demo_screenshot.png}
    \vspace{2em}
    \textit{[ Insert Live UI Screenshot Here ]} \\
    \vspace{1em}
    Demonstrating the resizable QSplitter layout with \\ the Waveform on top and the Spectrogram analyzer below.
  \end{center}
\end{frame}

\begin{frame}[standout]
  Thank You! \\
  \vspace{0.5em}
  \large Questions \& Answers
\end{frame}

\end{document}